%% ============================================================
%% VOLUME PROFISSIONAL — SONDAGEM E RASTREIO (SPEC-935-R210)
%% Integrante da colecao Alfabetizar Bem. Uso restrito a:
%% psicologo, psicopedagogo e neuropsicopedagogo.
%% NAO diagnostica; nao reproduz testes protegidos.
%% Faixa de referencia: 6-10 anos.
%% ============================================================
\documentclass[a4paper,11pt]{book}
\RequirePackage[brazil]{babel}
\RequirePackage{geometry}
\RequirePackage{booktabs}
\RequirePackage{longtable}
\RequirePackage{array}
\RequirePackage{enumitem}
\RequirePackage{xcolor}
\RequirePackage{fancyhdr}
\geometry{left=2.2cm, right=2.2cm, top=2.4cm, bottom=2.2cm}
\definecolor{cinza}{RGB}{90,90,90}
\pagestyle{fancy}
\fancyhf{}
\fancyhead[L]{\small\itshape Volume Profissional --- Alfabetizar Bem}
\fancyhead[R]{\small\thepage}
\renewcommand{\headrulewidth}{0.4pt}

\newcommand{\avisolegal}{\par\medskip\noindent\footnotesize\itshape
\color{cinza}Uso restrito a profissionais habilitados. Esta ficha NAO diagnostica;
serve ao rastreio e ao encaminhamento. Testes validados citados por referencia
(SATEPSI/CFP), sem reproducao. Resultados: sigilosos.\color{black}\par}

\begin{document}

\frontmatter
\begin{titlepage}
\begin{center}
\vspace*{1.5cm}
{\Huge\bfseries Volume Profissional}\\[0.7cm]
{\LARGE Sondagem e Rastreio}\\[0.5cm]
{\large Avaliacao do aluno e orientacao de professores}\\[1.2cm]
{\large Alfabetizar Bem --- Metodo Hibrido Fonico-Kumon}\\[2.2cm]
\noindent\rule{9cm}{0.6pt}\\[0.8cm]
{\normalsize\bfseries USO RESTRITO}\\[0.3cm]
{\small Este volume integra a colecao Alfabetizar Bem e destina-se a
psicologos, psicopedagogos e neuropsicopedagogos. Contem: Sondagem Inicial
minuciosa e completa (Parte A, antes da primeira atividade), fichas de
rastreio sistematico e investigativo de confirmacao (Parte B, 30 fichas) e
sintese/encaminhamento (Parte C).}\\[0.6cm]
{\small\itshape Nenhuma ficha deste volume gera diagnostico. Diagnostico e ato
privativo de profissional habilitado (Lei 4.119/1962; Res. CFP 009/2018).
Faixa de referencia: 6--10 anos.}
\end{center}
\end{titlepage}

\tableofcontents
\clearpage
\mainmatter

\chapter{Sondagem Inicial Minuciosa e Completa}

\section{Instrucoes de Aplicacao}

\begin{enumerate}[leftmargin=2.2em]
\item Aplicar \textbf{antes da primeira atividade} do volume correspondente.
\item Registrar para cada item: \textbf{0} = nao observado nesta faixa;
\textbf{1} = as vezes; \textbf{2} = frequentemente. Usar a coluna
\textit{Observacao} para exemplos objetivos.
\item Completar a leitura com: prova de escrita do nome, leitura de 5 palavras,
escrita de 5 palavras ditadas, copia de 3 frases curtas, contagem e
quantificacao ate 20, desenho da figura humana e nomeacao de 10 figuras.
\item Somar os pontos por dominio. \textbf{Prioridade de investigacao}:
dominio com maior soma. A soma NAO e nota, NAO e diagnostico e NAO substitui
teste validado; indica apenas onde aplicar as fichas de confirmacao (Parte B).
\item Sempre registrar data, idade e turma. Resultados sao sigilosos.
\end{enumerate}

\section{Perfil do Aluno}
\begin{tabular}{@{}ll@{}}
Nome: & \rule{8cm}{0.3pt}\\
Data da sondagem: & \rule{4cm}{0.3pt}\qquad Idade: \rule{2cm}{0.3pt}\qquad Turma: \rule{2cm}{0.3pt}\\
Aplicador (profissional): & \rule{7cm}{0.3pt}\\
Professor(a) colaborador(a): & \rule{6.5cm}{0.3pt}\\
\end{tabular}
\medskip

\section{Dominio 1 --- Leitura}
\begin{quote}\itshape Escala: 0 = nao observado; 1 = as vezes; 2 = frequentemente.\end{quote}
\begin{tabular}{@{}p{10.5cm}ccc p{0.5cm}@{}}\toprule
Item & 0 & 1 & 2 & Obs. \\\midrule
Reconhece letras do proprio nome & & & & \\
Nomeia letras fora de ordem & & & & \\
Le silabas simples (ba, la, me) & & & & \\
Le palavras com estrutura CVC (pato, bolo) & & & & \\
Le pseudopalavras sem ajuda de adivinhacao & & & & \\
Segmenta palavras em silabas ao ler & & & & \\
Compreende o que le (reconta ideia central) & & & & \\
Troca/omite letras ao ler (ex.: pate por pato) & & & & \\
Le devagar, soletrando & & & & \\
Evita situacoes de leitura & & & & \\
\bottomrule\end{tabular}
\vspace{0.2cm}\noindent\small Total do dominio: \rule{1cm}{0.3pt} / 20\par\medskip


\section{Dominio 2 --- Escrita / Ortografia}
\begin{quote}\itshape Escala: 0 = nao observado; 1 = as vezes; 2 = frequentemente.\end{quote}
\begin{tabular}{@{}p{10.5cm}ccc p{0.5cm}@{}}\toprule
Item & 0 & 1 & 2 & Obs. \\\midrule
Escreve o proprio nome sem copiar & & & & \\
Copia textos sem erros de transposicao & & & & \\
Escreve palavras ditadas com estrutura simples & & & & \\
Representa todos os fonemas ao escrever & & & & \\
Respeita espacamento entre palavras & & & & \\
Usa maiuscula no inicio e ponto final & & & & \\
Erros ortograficos regulares persistentes & & & & \\
Confunde letras de traco similar (b/d, p/q) & & & & \\
Omissao de letras em palavras (fpa por foca) & & & & \\
Fadiga ou dor ao escrever & & & & \\
\bottomrule\end{tabular}
\vspace{0.2cm}\noindent\small Total do dominio: \rule{1cm}{0.3pt} / 20\par\medskip


\section{Dominio 3 --- Matematica}
\begin{quote}\itshape Escala: 0 = nao observado; 1 = as vezes; 2 = frequentemente.\end{quote}
\begin{tabular}{@{}p{10.5cm}ccc p{0.5cm}@{}}\toprule
Item & 0 & 1 & 2 & Obs. \\\midrule
Conta ate 20 de forma estavel & & & & \\
Relaciona numero a quantidade ate 10 & & & & \\
Compara quantidades (mais/menos) & & & & \\
Soma e subtrai ate 10 com material concreto & & & & \\
Reconhece numeros ate 20 fora de ordem & & & & \\
Copla padroes e sequencias & & & & \\
Compreende dezena/unidade & & & & \\
Resolve problemas orais simples & & & & \\
Confunde sinais + e - & & & & \\
Evita atividades numericas & & & & \\
\bottomrule\end{tabular}
\vspace{0.2cm}\noindent\small Total do dominio: \rule{1cm}{0.3pt} / 20\par\medskip


\section{Dominio 4 --- Consciencia Fonologica}
\begin{quote}\itshape Escala: 0 = nao observado; 1 = as vezes; 2 = frequentemente.\end{quote}
\begin{tabular}{@{}p{10.5cm}ccc p{0.5cm}@{}}\toprule
Item & 0 & 1 & 2 & Obs. \\\midrule
Rima palavras (pato/gato) & & & & \\
Segmenta palavras em silabas & & & & \\
Identifica silaba inicial & & & & \\
Identifica silaba final & & & & \\
Aliteracao (mesmo som inicial) & & & & \\
Sintetiza fonemas (f-a-l-a = fala) & & & & \\
Segmenta fonemas em palavras simples & & & & \\
Manipula fonemas (troca /p/ por /t/ em pato) & & & & \\
Percebe sons iguais em palavras diferentes & & & & \\
Memoriza sequencias orais (dias, rimas) & & & & \\
\bottomrule\end{tabular}
\vspace{0.2cm}\noindent\small Total do dominio: \rule{1cm}{0.3pt} / 20\par\medskip


\section{Dominio 5 --- Atencao e Funcoes Executivas}
\begin{quote}\itshape Escala: 0 = nao observado; 1 = as vezes; 2 = frequentemente.\end{quote}
\begin{tabular}{@{}p{10.5cm}ccc p{0.5cm}@{}}\toprule
Item & 0 & 1 & 2 & Obs. \\\midrule
Mantem atencao em tarefa de 10 minutos & & & & \\
Atende instrucoes orais de 2 passos & & & & \\
Organiza material e mochila & & & & \\
Inicia tarefa sem adiamento excessivo & & & & \\
Persiste diante de erro sem desistir & & & & \\
Controla impulsos (espera vez) & & & & \\
Planeja passos de uma atividade & & & & \\
Autorregula ritmo de trabalho & & & & \\
Conclui tarefas iniciadas & & & & \\
Alterna entre atividades com apoio minimo & & & & \\
\bottomrule\end{tabular}
\vspace{0.2cm}\noindent\small Total do dominio: \rule{1cm}{0.3pt} / 20\par\medskip


\section{Dominio 6 --- Linguagem Oral}
\begin{quote}\itshape Escala: 0 = nao observado; 1 = as vezes; 2 = frequentemente.\end{quote}
\begin{tabular}{@{}p{10.5cm}ccc p{0.5cm}@{}}\toprule
Item & 0 & 1 & 2 & Obs. \\\midrule
Compreende ordens simples e complexas & & & & \\
Vocabulario adequado a idade & & & & \\
Narra acontecimentos em sequencia & & & & \\
Usa frases completas & & & & \\
Encontra palavras (acesso lexical) & & & & \\
Pronuncia todos os fonemas da lingua & & & & \\
Compreende perguntas quem/onde/por que & & & & \\
Conta historias com coesao & & & & \\
Participa de dialogo mantendo o tema & & & & \\
Reconta historias ouvidas com detalhes & & & & \\
\bottomrule\end{tabular}
\vspace{0.2cm}\noindent\small Total do dominio: \rule{1cm}{0.3pt} / 20\par\medskip


\section{Dominio 7 --- Motricidade Fina / Grafomotricidade}
\begin{quote}\itshape Escala: 0 = nao observado; 1 = as vezes; 2 = frequentemente.\end{quote}
\begin{tabular}{@{}p{10.5cm}ccc p{0.5cm}@{}}\toprule
Item & 0 & 1 & 2 & Obs. \\\midrule
Segura lápis com preensao funcional & & & & \\
Recorta com tesoura & & & & \\
Desenha figura humana com partes essenciais & & & & \\
Copla formas geometricas & & & & \\
Escreve com pressao adequada & & & & \\
Respeita o espaco da linha ao escrever & & & & \\
Conclui tracados sem inversao exagerada & & & & \\
Coordena olho-mao em labirintos & & & & \\
Distingue direita/esquerda no proprio corpo & & & & \\
Apresenta movimentos lentos ou imprecisos & & & & \\
\bottomrule\end{tabular}
\vspace{0.2cm}\noindent\small Total do dominio: \rule{1cm}{0.3pt} / 20\par\medskip


\section{Dominio 8 --- Comportamento e Socioemocional}
\begin{quote}\itshape Escala: 0 = nao observado; 1 = as vezes; 2 = frequentemente.\end{quote}
\begin{tabular}{@{}p{10.5cm}ccc p{0.5cm}@{}}\toprule
Item & 0 & 1 & 2 & Obs. \\\midrule
Interage com pares sem conflitos frequentes & & & & \\
Aceita regras e limites & & & & \\
Expressa emocoes com palavras & & & & \\
Busca adultos quando necessario & & & & \\
Tolera mudancas de rotina & & & & \\
Participa de brincadeiras coletivas & & & & \\
Reage com intensidade desproporcional & & & & \\
Fica retraido ou isolado & & & & \\
Apresenta medos intensos de separacao & & & & \\
Apresenta preocupacoes repetitivas & & & & \\
\bottomrule\end{tabular}
\vspace{0.2cm}\noindent\small Total do dominio: \rule{1cm}{0.3pt} / 20\par\medskip


\section{Dominio 9 --- Processamento Auditivo}
\begin{quote}\itshape Escala: 0 = nao observado; 1 = as vezes; 2 = frequentemente.\end{quote}
\begin{tabular}{@{}p{10.5cm}ccc p{0.5cm}@{}}\toprule
Item & 0 & 1 & 2 & Obs. \\\midrule
Compreende em ambiente ruidoso & & & & \\
Repete sequencias de 3-4 palavras & & & & \\
Distingue sons semelhantes da fala & & & & \\
Segue instrucoes longas sem repeticao & & & & \\
Memoriza musicas e rimas & & & & \\
Apresenta pedidos frequentes de repeticao & & & & \\
Confunde palavras parecidas ao ouvir & & & & \\
Percebe variacoes de entonacao & & & & \\
Utiliza pistas visuais para compensar & & & & \\
Demonstra cansaco ao ouvir explicacoes & & & & \\
\bottomrule\end{tabular}
\vspace{0.2cm}\noindent\small Total do dominio: \rule{1cm}{0.3pt} / 20\par\medskip


\section{Dominio 10 --- Processamento Visual / Visomotor}
\begin{quote}\itshape Escala: 0 = nao observado; 1 = as vezes; 2 = frequentemente.\end{quote}
\begin{tabular}{@{}p{10.5cm}ccc p{0.5cm}@{}}\toprule
Item & 0 & 1 & 2 & Obs. \\\midrule
Copia do quadro sem erros de posicao & & & & \\
Completa labirintos sem sair da trilha & & & & \\
Distingue letras e formas parecidas & & & & \\
Mantem orientacao no papel (nao inverte) & & & & \\
Percebe diferencas em figuras & & & & \\
Monta quebra-cabecas adequados a idade & & & & \\
Copia figuras com proporcao & & & & \\
Localiza itens em cena complexa & & & & \\
Escreve respeitando direcao da linha & & & & \\
Apresenta queixas visuais (dore de cabeca, aperta olhos) & & & & \\
\bottomrule\end{tabular}
\vspace{0.2cm}\noindent\small Total do dominio: \rule{1cm}{0.3pt} / 20\par\medskip


\section{Dominio 11 --- Saude Sensorial (sinais para encaminhamento)}
\begin{quote}\itshape Escala: 0 = nao observado; 1 = as vezes; 2 = frequentemente.\end{quote}
\begin{tabular}{@{}p{10.5cm}ccc p{0.5cm}@{}}\toprule
Item & 0 & 1 & 2 & Obs. \\\midrule
Reage a chamado em volume normal & & & & \\
Percebe sons do ambiente & & & & \\
Acompanha sons para a fonte & & & & \\
Realiza leitura com distancia adequada & & & & \\
Espreme os olhos para ver o quadro & & & & \\
Apresenta baixa resposta a sons agudos & & & & \\
Apresenta historico de otites ou perdas & & & & \\
Usa dispositivo assistivo declarado & & & & \\
Troca letras por dificuldade auditiva & & & & \\
Apresenta fonemas mal articulados por audicao & & & & \\
\bottomrule\end{tabular}
\vspace{0.2cm}\noindent\small Total do dominio: \rule{1cm}{0.3pt} / 20\par\medskip


\clearpage

\chapter*{Parte B --- Fichas de Rastreio Sistematico e Investigativo (Confirmacao)}
\addcontentsline{toc}{chapter}{Parte B --- Fichas de Rastreio}

\chapter{Dislexia}
\section*{Dominio: Leitura}
\section*{Ficha de rastreio sistematico -- confirmacao}
\subsection*{Sinais de alerta (marque: S = sim, AV = as vezes, N = nao)}
\begin{tabular}{@{}p{11.5cm}ccc@{}}\toprule
Sinal observado & S & AV & N \\\midrule
Leitura lenta, soletrada, com muitas pausas & & & \\
Erros de inversao (fada/dafa), omissao e adicao & & & \\
Dificuldade de rima e aliteracao persistente & & & \\
Nao reconhece palavras familiares de forma automatica & & & \\
Compreensao prejudicada pelo esforco de decodificacao & & & \\
Historico familiar de dificuldade de leitura & & & \\
Evita ler em voz alta & & & \\
Desempenho muito abaixo do esperado para a idade/instrucao & & & \\
Melhora quando o texto e lido por outra pessoa & & & \\
Fadiga rapida em tarefas de leitura & & & \\
\bottomrule\end{tabular}
\subsection*{Como investigar (tarefas sistematicas)}
\begin{itemize}
\item Pedir leitura de 5 palavras e 3 pseudopalavras com cronometro
\item Rima e segmentacao fonemica oral
\item Leitura de texto em voz alta com registro de erros
\item Comparar com producao de leitura de mesma turma
\end{itemize}
\subsection*{Exclusoes e confundidores}
\begin{itemize}
\item Falta de exposicao a leitura
\item Problema visual nao corrigido
\item Ansiedade de desempenho pontual
\item Ensino inicial inconsistente (metodos exclusivamente nao-fonicos)
\end{itemize}
\subsection*{Instrumentos validados para confirmacao (referencia -- aplicar somente por psicologo, conforme SATEPSI; nao reproduzidos aqui)}
\begin{itemize}
\item TDE-II (subtestes de leitura) -- SATEPSI
\item PROLEC / PROOHFON (perfil fonologico)
\item Avaliacao neuropsicologica (psicologo)
\end{itemize}
\subsection*{Encaminhamento}
\begin{itemize}
\item Psicologo/neuropsicologo para avaliacao formal
\item Fonoaudiologo para consciencia fonologica
\item Pedagogo especializado se confirmado
\end{itemize}
\subsection*{Adaptacoes pedagogicas iniciais (sem rotulo)}
\begin{itemize}
\item Leitura compartilhada diaria com apoio fonico
\item Letra ampliada e textos curtos
\item Ensinar rota fonologica explicita (fonema-grafema)
\item Prova oral como alternativa a prova escrita
\item Tempo adicional nas atividades
\end{itemize}
\subsection*{Bibliografia basica}
\begin{itemize}
\item Shaywitz (2006), Dislexia. Artmed.
\item Capovilla \& Capovilla (2007), Consciencia Fonologica. Vetor.
\end{itemize}
\avisolegal


\chapter{Disortografia}
\section*{Dominio: Escrita/Ortografia}
\section*{Ficha de rastreio sistematico -- confirmacao}
\subsection*{Sinais de alerta (marque: S = sim, AV = as vezes, N = nao)}
\begin{tabular}{@{}p{11.5cm}ccc@{}}\toprule
Sinal observado & S & AV & N \\\midrule
Erros ortograficos muito alem do esperado para a serie & & & \\
Omissao, troca e adicao de letras em palavras regulares & & & \\
Inversoes e transposicoes (pato/vato) & & & \\
Dificuldade com regras contextuais (R/RR, S/SS) & & & \\
Nao usa acentos e marcadores de nasalidade & & & \\
Dificuldade com palavras irregulares de alta frequencia & & & \\
Escrita lenta e laboriosa & & & \\
Contraste entre ideia verbal rica e texto escrito pobre & & & \\
\bottomrule\end{tabular}
\subsection*{Como investigar (tarefas sistematicas)}
\begin{itemize}
\item Ditado de 20 palavras regulares e irregulares
\item Escrita espontanea (tema curto) e analise de erros
\item Prova de conhecimento das regras ortograficas
\item Comparar com escrita de pares da mesma turma
\end{itemize}
\subsection*{Exclusoes e confundidores}
\begin{itemize}
\item Pouca experiencia de escrita
\item Erros de transposicao visual (ver ficha 3)
\item Problema de atencao sustentada
\end{itemize}
\subsection*{Instrumentos validados para confirmacao (referencia -- aplicar somente por psicologo, conforme SATEPSI; nao reproduzidos aqui)}
\begin{itemize}
\item TDE-II (subteste de escrita)
\item Avaliacao de linguagem escrita (psicopedagogo)
\item Perfil ortografico (fonoaudioloogia)
\end{itemize}
\subsection*{Encaminhamento}
\begin{itemize}
\item Psicopedagogo para reeducacao da escrita
\item Fonoaudiologo se houver trocas de fonemas
\item Psicologo se associado a TDAH/dislexia
\end{itemize}
\subsection*{Adaptacoes pedagogicas iniciais (sem rotulo)}
\begin{itemize}
\item Reforco das regras contextuais com jogos
\item Revisao dialogada do texto (o professor pergunta, aluno encontra)
\item Lista pessoal de palavras-alvo
\item Evitar castigo por erros; focar em 1 padrao por semana
\item Uso de corretor ortografico assistido no computador
\end{itemize}
\subsection*{Bibliografia basica}
\begin{itemize}
\item Zorzi (1998), Aprender a escrever. Artmed.
\item Scliar-Cabral (2003), Guia pratico de alfabetizacao. Contexto.
\end{itemize}
\avisolegal


\chapter{Disgrafia}
\section*{Dominio: Motricidade Fina/Grafomotricidade}
\section*{Ficha de rastreio sistematico -- confirmacao}
\subsection*{Sinais de alerta (marque: S = sim, AV = as vezes, N = nao)}
\begin{tabular}{@{}p{11.5cm}ccc@{}}\toprule
Sinal observado & S & AV & N \\\midrule
Caligrafia ilegivel ou muito irregular & & & \\
Preensao inadequada do lapis com tensao & & & \\
Formatos de letras inconsistentes & & & \\
Distancias entre letras e palavras desiguais & & & \\
Escrita muito lenta e dolorosa & & & \\
Postura inadequada ao escrever & & & \\
Coordenacao fina pobre (abotoar, recortar) & & & \\
Diferenca grande entre texto oral e escrito & & & \\
\bottomrule\end{tabular}
\subsection*{Como investigar (tarefas sistematicas)}
\begin{itemize}
\item Observar preensao, postura e velocidade
\item Copia de frases com registro de legibilidade
\item Teste de destreza manual (recorte, enfiar contas)
\item Excluir dor/fadiga durante e apos escrita
\end{itemize}
\subsection*{Exclusoes e confundidores}
\begin{itemize}
\item Pressa ou motivacao
\item Falta de treino de escrita
\item Problema visual
\item Dificuldade de consciencia fonologica (disortografia)
\end{itemize}
\subsection*{Instrumentos validados para confirmacao (referencia -- aplicar somente por psicologo, conforme SATEPSI; nao reproduzidos aqui)}
\begin{itemize}
\item Avaliacao psicomotora (terapeuta ocupacional/psicomotricista)
\item Avaliacao neuropsicologica se necessario
\end{itemize}
\subsection*{Encaminhamento}
\begin{itemize}
\item Terapeuta ocupacional para reeducacao grafomotora
\item Psicopedagogo para estrategias de escrita
\item Ortopedista se dor persistente
\end{itemize}
\subsection*{Adaptacoes pedagogicas iniciais (sem rotulo)}
\begin{itemize}
\item Posicao do papel inclinada e apoio de antebraco
\item Instrumentos de escrita ergonomicos (lapis triangular)
\item Pauta com linhas de apoio (pauta dupla, Mestra4)
\item Tempo reduzido de copia, com foco em qualidade
\item Uso de teclado/tablet como alternativa
\end{itemize}
\subsection*{Bibliografia basica}
\begin{itemize}
\item Fonseca (2012), Manual de avaliacao psicomotora. Ancora.
\item Oliveira (2010), Manual de psicomotricidade. Wak.
\end{itemize}
\avisolegal


\chapter{Discalculia}
\section*{Dominio: Matematica}
\section*{Ficha de rastreio sistematico -- confirmacao}
\subsection*{Sinais de alerta (marque: S = sim, AV = as vezes, N = nao)}
\begin{tabular}{@{}p{11.5cm}ccc@{}}\toprule
Sinal observado & S & AV & N \\\midrule
Dificuldade persistente em contar de forma estavel & & & \\
Nao reconhece numeros fora de ordem & & & \\
Incapacidade de relacionar numero a quantidade alem do concreto & & & \\
Erros com valor posicional (dezena/unidade) & & & \\
Dificuldade com fatos basicos (2+3) sem contagem manual & & & \\
Problemas com comparacao (mais/menos) e estimativa & & & \\
Confusao de sinais e procedimentos apesar de pratica & & & \\
Erros de calculo apesar de dominar o procedimento com apoio em calculo e raciocinio & & & \\
Evita numeros e jogos de conta & & & \\
Dificuldade com relogio, dinheiro, medidas & & & \\
\bottomrule\end{tabular}
\subsection*{Como investigar (tarefas sistematicas)}
\begin{itemize}
\item Prova de contagem, transcodificacao e calculo mental
\item Tarefas de comparacao de quantidades (estimativa)
\item Resolucao de problemas com apoio concreto
\item Observar uso de estrategias (contar nos dedos vs. calculo mental)
\end{itemize}
\subsection*{Exclusoes e confundidores}
\begin{itemize}
\item Ansiedade matematica
\item Pouca exposicao a numeros
\item Dificuldade atencional (TDAH)
\item Problema de linguagem (interpretacao)
\end{itemize}
\subsection*{Instrumentos validados para confirmacao (referencia -- aplicar somente por psicologo, conforme SATEPSI; nao reproduzidos aqui)}
\begin{itemize}
\item Avaliacao neuropsicologica (psicologo)
\item Testes de desempenho em matematica (TDE-II)
\item Raciocinio quantitativo (WISC-V)
\end{itemize}
\subsection*{Encaminhamento}
\begin{itemize}
\item Psicopedagogo para reeducacao matematica
\item Psicologo/neuropsicologo para avaliacao formal
\end{itemize}
\subsection*{Adaptacoes pedagogicas iniciais (sem rotulo)}
\begin{itemize}
\item Usar material concreto (material dourado, fichas)
\item Ensinar uma estrategia por vez com sobrepratica
\item Quadro de numeros e reta numerica visiveis
\item Problemas orais antes de escritos
\item Minimizar copia de enunciados longos
\end{itemize}
\subsection*{Bibliografia basica}
\begin{itemize}
\item Butterworth \& Yeo (2009), Discalculia. Bookman.
\item Kaufmann et al. (2010), Desenvolvimento do calculo. Artmed.
\end{itemize}
\avisolegal


\chapter{TDAH -- desatencao}
\section*{Dominio: Atencao e Funcoes Executivas}
\section*{Ficha de rastreio sistematico -- confirmacao}
\subsection*{Sinais de alerta (marque: S = sim, AV = as vezes, N = nao)}
\begin{tabular}{@{}p{11.5cm}ccc@{}}\toprule
Sinal observado & S & AV & N \\\midrule
Dificuldade de manter atencao em detalhes, erros por descuido & & & \\
Nao parece ouvir quando falam diretamente & & & \\
Nao segue instrucoes e nao conclui tarefas & & & \\
Dificuldade de organizar material e tempo & & & \\
Evita/esforco com tarefas que exigem atencao sustentada & & & \\
Perde objetos (material, agendas) & & & \\
Distrai-se facilmente com estimulos externos & & & \\
Esquecimento nas atividades diarias & & & \\
Sinais presentes em casa e na escola & & & \\
Inicio antes dos 12 anos & & & \\
\bottomrule\end{tabular}
\subsection*{Como investigar (tarefas sistematicas)}
\begin{itemize}
\item Escalas comportamentais respondidas por pais e professores
\item Observacao sistematica em tarefa de 10 minutos
\item Entrevista de desenvolvimento e historico escolar
\item Registrar se os sinais ocorrem em dois contextos
\end{itemize}
\subsection*{Exclusoes e confundidores}
\begin{itemize}
\item Ansiedade (preocupacoes explicam o desatento)
\item Ambiente muito estimulante
\item Sono inadequado
\item Perda auditiva discreta
\end{itemize}
\subsection*{Instrumentos validados para confirmacao (referencia -- aplicar somente por psicologo, conforme SATEPSI; nao reproduzidos aqui)}
\begin{itemize}
\item Escalas SNAP-IV / ASRS (triagem; uso profissional)
\item WISC-V / testes de atencao (psicologo, SATEPSI)
\item Avaliacao psiquiatrica se necessario
\end{itemize}
\subsection*{Encaminhamento}
\begin{itemize}
\item Psicologo/neuropsicologo para avaliacao
\item Pediatra/psiquiatra infantil se confirmado
\item Psicopedagogo para organizacao de estudos
\end{itemize}
\subsection*{Adaptacoes pedagogicas iniciais (sem rotulo)}
\begin{itemize}
\item Cadeira proxima ao professor e longe de fontes de distracao
\item Instrucoes curtas e fracionadas (1 passo por vez)
\item Checklist visual de rotina na carteira
\item Pausas motoras programadas (5 min)
\item Reforco imediato de comportamento desejado
\end{itemize}
\subsection*{Bibliografia basica}
\begin{itemize}
\item Barkley (2008), Transtorno de deficit de atencao. Artmed.
\item Rohde \& Mattos (2003), Principios e praticas em TDAH. Artmed.
\end{itemize}
\avisolegal


\chapter{TDAH -- hiperatividade/impulsividade}
\section*{Dominio: Atencao e Comportamento}
\section*{Ficha de rastreio sistematico -- confirmacao}
\subsection*{Sinais de alerta (marque: S = sim, AV = as vezes, N = nao)}
\begin{tabular}{@{}p{11.5cm}ccc@{}}\toprule
Sinal observado & S & AV & N \\\midrule
Remexe-se, batate as maos ou pes & & & \\
Levanta-se quando deveria permanecer sentado & & & \\
Corre/escala em situacoes inapropriadas & & & \\
Fala excessivamente e nao brinca em silencio & & & \\
Responde antes de a pergunta terminar & & & \\
Dificuldade de esperar a vez em filas/jogos & & & \\
Interrompe intromete-se nos outros & & & \\
Age sem avaliar consequencias (acidentes) & & & \\
Apresenta os sinais em mais de um contexto & & & \\
Deficit funcional em casa e na escola & & & \\
\bottomrule\end{tabular}
\subsection*{Como investigar (tarefas sistematicas)}
\begin{itemize}
\item Escalas SNAP-IV / ASRS (triagem)
\item Observacao em fila, roda e atividade livre
\item Registro de episodios com antecedentes e consequencias
\end{itemize}
\subsection*{Exclusoes e confundidores}
\begin{itemize}
\item Ansiedade/trauma que aumenta agitacao
\item Sono insuficiente
\item Temperamento forte sem prejuizo funcional
\end{itemize}
\subsection*{Instrumentos validados para confirmacao (referencia -- aplicar somente por psicologo, conforme SATEPSI; nao reproduzidos aqui)}
\begin{itemize}
\item Escalas e entrevista clinica (psicologo)
\item Avaliacao psiquiatrica se prejuizo grave
\end{itemize}
\subsection*{Encaminhamento}
\begin{itemize}
\item Psicologo/psiquiatra infantil
\item Terapia comportamental parental se indicado
\end{itemize}
\subsection*{Adaptacoes pedagogicas iniciais (sem rotulo)}
\begin{itemize}
\item Combinados claros e visuais com consequencias previsiveis
\item Pausas motoras regulares
\item Lugar estrategico na sala
\item Ensinar a parar-pensar-executar com cartao
\item Reduzir estimulos concorrentes em provas
\end{itemize}
\subsection*{Bibliografia basica}
\begin{itemize}
\item Barkley (2008), Transtorno de deficit de atencao. Artmed.
\item DSM-5-TR (2023). Artmed.
\end{itemize}
\avisolegal


\chapter{TEA -- nivel 1}
\section*{Dominio: Comportamento/Social}
\section*{Ficha de rastreio sistematico -- confirmacao}
\subsection*{Sinais de alerta (marque: S = sim, AV = as vezes, N = nao)}
\begin{tabular}{@{}p{11.5cm}ccc@{}}\toprule
Sinal observado & S & AV & N \\\midrule
Dificuldade de reciprocidade social (nao inicia/retribui interacoes) & & & \\
Interesses restritos e intensos (trens, espacos, letras) & & & \\
Rotinas inflexiveis com reacao intensa a mudancas & & & \\
Dificuldade de leitura de expressoes e ironia & & & \\
Fala fluente porem com dificuldade pragmática & & & \\
Pouco contato visual funcional & & & \\
Dificuldade de fazer e manter amigos da mesma idade & & & \\
Hipersensibilidade ou hipossensibilidade sensorial & & & \\
Dificuldade em brincar de faz de conta com pares & & & \\
Sinais presentes desde a infancia & & & \\
\bottomrule\end{tabular}
\subsection*{Como investigar (tarefas sistematicas)}
\begin{itemize}
\item Observacao estruturada em recreio e atividades livres
\item Entrevista de desenvolvimento (marcos sociais/ludio)
\item Prova de compreensao pragmatica (ironias, figurativa)
\end{itemize}
\subsection*{Exclusoes e confundidores}
\begin{itemize}
\item Timidez/ansiedade social
\item TDL (linguagem) sem compromisso social
\item Perfil superdotado com interesses intensos
\end{itemize}
\subsection*{Instrumentos validados para confirmacao (referencia -- aplicar somente por psicologo, conforme SATEPSI; nao reproduzidos aqui)}
\begin{itemize}
\item ADOS-2 (referencia internacional; validação brasileira em andamento) / ADI-R (validado no Brasil; Becker et al., 2012)
\item Escala SRS-2 (triagem; uso profissional)
\item Avaliacao neuropsicologica
\end{itemize}
\subsection*{Encaminhamento}
\begin{itemize}
\item Psicologo/neuropsicologo para avaliacao formal
\item Fonoaudiologo para pragmatica e suporte
\item Psiquiatra infantil se necessario
\end{itemize}
\subsection*{Adaptacoes pedagogicas iniciais (sem rotulo)}
\begin{itemize}
\item Rotina previsivel com aviso de transicoes
\item Suportes visuais (agenda, regras)
\item Grupo de habilidades sociais
\item Ensinar explicitamente trocas sociais
\item Ajuste sensorial (fones, espaco tranquilo)
\end{itemize}
\subsection*{Bibliografia basica}
\begin{itemize}
\item APA (2023), DSM-5-TR. Artmed.
\item Klin, Volkmar \& Sparrow (2000). Asperger syndrome. Guilford.
\end{itemize}
\avisolegal


\chapter{TEA -- nivel 2/3}
\section*{Dominio: Comportamento/Social/Linguagem}
\section*{Ficha de rastreio sistematico -- confirmacao}
\subsection*{Sinais de alerta (marque: S = sim, AV = as vezes, N = nao)}
\begin{tabular}{@{}p{11.5cm}ccc@{}}\toprule
Sinal observado & S & AV & N \\\midrule
Atraso marcado ou ausencia de linguagem funcional & & & \\
Comunicacao muito reduzida mesmo com gestos alternativos & & & \\
Respostas sociais ausentes ou minimas & & & \\
Interesses e comportamentos repetitivos intensos & & & \\
Rituais e resistencia intensa a mudanca & & & \\
Autoagressao ou agressao quando contrariado (sinais de comunicar) & & & \\
Pouco uso de olhar, gestos e expressao para comunicar & & & \\
Bom progresso com suporte visual intensivo & & & \\
Ecolalia persistente ale da idade esperada & & & \\
Comprometimento severo na participacao escolar & & & \\
\bottomrule\end{tabular}
\subsection*{Como investigar (tarefas sistematicas)}
\begin{itemize}
\item Observacao em rotina estruturada
\item Entrevista com responsaveis (marcos)
\item Avaliacao de comunicacao funcional (pedido, recusa)
\end{itemize}
\subsection*{Exclusoes e confundidores}
\begin{itemize}
\item Deficiencia auditiva nao tratada
\item DI moderada sem prejuizo social especifico do TEA
\end{itemize}
\subsection*{Instrumentos validados para confirmacao (referencia -- aplicar somente por psicologo, conforme SATEPSI; nao reproduzidos aqui)}
\begin{itemize}
\item ADOS-2 (referencia internacional; validação brasileira em andamento) / ADI-R (validado no Brasil; Becker et al., 2012)
\item Avaliacao de linguagem (fonoaudiologo)
\item Avaliacao medica/neurologica
\end{itemize}
\subsection*{Encaminhamento}
\begin{itemize}
\item Neurologista/psiquiatra infantil + psicologo
\item Fonoaudiologo intensivo
\item Terapeuta ocupacional/ABA se disponivel
\end{itemize}
\subsection*{Adaptacoes pedagogicas iniciais (sem rotulo)}
\begin{itemize}
\item Comunicacao alternativa (PECS/gestos) se necessaria
\item Rotina visual clara com aviso de mudanca
\item Espaco sensorial calmo
\item Fracionar atividades com reforco
\item Mentor/mediador de apoio em transicoes
\end{itemize}
\subsection*{Bibliografia basica}
\begin{itemize}
\item APA (2023), DSM-5-TR. Artmed.
\item Suzuki \& Barros (2009). TEA na escola. Vetor.
\end{itemize}
\avisolegal


\chapter{Transtorno Opositivo-Desafiador (TOD)}
\section*{Dominio: Comportamento}
\section*{Ficha de rastreio sistematico -- confirmacao}
\subsection*{Sinais de alerta (marque: S = sim, AV = as vezes, N = nao)}
\begin{tabular}{@{}p{11.5cm}ccc@{}}\toprule
Sinal observado & S & AV & N \\\midrule
Perde a calma com frequencia & & & \\
Discute com adultos e desafia regras & & & \\
Recusa cumprir pedidos de autoridades & & & \\
Irrita deliberadamente pessoas & & & \\
Culpa os outros pelos proprios erros & & & \\
Ressentido ou vingativo & & & \\
Sinais presentes ha pelo menos 6 meses & & & \\
Prejuizo nas relacoes com pares e professores & & & \\
Comportamentos nao explicados por humor depressivo & & & \\
Padrao consistente em casa e/ou escola & & & \\
\bottomrule\end{tabular}
\subsection*{Como investigar (tarefas sistematicas)}
\begin{itemize}
\item Registro de episodios (gatilhos, duracao)
\item Escala de comportamento (CBCL/TRF -- referencia)
\item Entrevista com pais sobre limites e consequencias
\end{itemize}
\subsection*{Exclusoes e confundidores}
\begin{itemize}
\item TDAH com impulsividade
\item Depressao/irritabilidade
\item Estilo parental inconsistente
\item Bullying/ambiente hostil
\end{itemize}
\subsection*{Instrumentos validados para confirmacao (referencia -- aplicar somente por psicologo, conforme SATEPSI; nao reproduzidos aqui)}
\begin{itemize}
\item CBCL/TRF (uso profissional)
\item Entrevista diagnostica (psicologo/psiquiatra)
\end{itemize}
\subsection*{Encaminhamento}
\begin{itemize}
\item Psicologo (terapia comportamental/familiar)
\item Psiquiatra infantil se comorbidades
\end{itemize}
\subsection*{Adaptacoes pedagogicas iniciais (sem rotulo)}
\begin{itemize}
\item Reforco positivo frequente (4:1)
\item Consequencias previsiveis e imediatas
\item Escolhas limitadas (isto ou aquilo)
\item Ensinar e reforcar habilidades de resolucao de problemas
\item Comunicacao com a familia alinhada
\end{itemize}
\subsection*{Bibliografia basica}
\begin{itemize}
\item DSM-5-TR (2023). Artmed.
\item Weber (2017), Comportamento infantil. Artmed.
\end{itemize}
\avisolegal


\chapter{Transtorno de Conduta}
\section*{Dominio: Comportamento}
\section*{Ficha de rastreio sistematico -- confirmacao}
\subsection*{Sinais de alerta (marque: S = sim, AV = as vezes, N = nao)}
\begin{tabular}{@{}p{11.5cm}ccc@{}}\toprule
Sinal observado & S & AV & N \\\midrule
Agressao fisica a pares/pessoas & & & \\
Mentiras repetidas para obter vantagem & & & \\
Vandalismo ou destruicao de bens & & & \\
Furtos (inclusive na escola) & & & \\
Violacao grave de regras (fugas) & & & \\
Crueldade com animais & & & \\
Inicio antes dos 13 anos em alguns casos & & & \\
Falta de remorso apos dano & & & \\
Padrao persistente por 12 meses & & & \\
Prejuizo social e academico significativo & & & \\
\bottomrule\end{tabular}
\subsection*{Como investigar (tarefas sistematicas)}
\begin{itemize}
\item Registro sistematico de incidentes
\item Entrevista com familia e escola
\item Avaliacao de contexto (violencia domestica, grupo)
\end{itemize}
\subsection*{Exclusoes e confundidores}
\begin{itemize}
\item TDAH impulsivo
\item TOD (sem violacao de direitos)
\item Contexto de violencia/negligencia (nao diagnosticar o contexto como transtorno)
\end{itemize}
\subsection*{Instrumentos validados para confirmacao (referencia -- aplicar somente por psicologo, conforme SATEPSI; nao reproduzidos aqui)}
\begin{itemize}
\item Avaliacao clinica/psiquiatrica
\item Instrumentos de comportamento (CBCL/TRF)
\end{itemize}
\subsection*{Encaminhamento}
\begin{itemize}
\item Psicologo/psiquiatra infantil
\item Rede de protecao (Conselho Tutelar, CRAS) se necessario
\item Escola: plano de convivencia
\end{itemize}
\subsection*{Adaptacoes pedagogicas iniciais (sem rotulo)}
\begin{itemize}
\item Limites claros e consistentes com reparacao de dano
\item Ensino de empatia por historias e role-play
\item Supervisao positiva e rotina estruturada
\item Contracto de comportamento com metas
\item Articulacao com rede de protecao
\end{itemize}
\subsection*{Bibliografia basica}
\begin{itemize}
\item DSM-5-TR (2023). Artmed.
\item Patterson (2002), Antisocial boys. Castalia.
\end{itemize}
\avisolegal


\chapter{Ansiedade de Separacao}
\section*{Dominio: Socioemocional}
\section*{Ficha de rastreio sistematico -- confirmacao}
\subsection*{Sinais de alerta (marque: S = sim, AV = as vezes, N = nao)}
\begin{tabular}{@{}p{11.5cm}ccc@{}}\toprule
Sinal observado & S & AV & N \\\midrule
Medo intenso de separar-se dos pais & & & \\
Choro ou crises ao entrar na escola & & & \\
Sintomas somaticos (dor de barriga, enjoo) em dias de aula & & & \\
Preocupacao de que algo aconteca aos pais & & & \\
Pesadelos sobre separacao & & & \\
Recusa de dormir fora de casa & & & \\
Segue os pais pela casa & & & \\
Sintomas ha pelo menos 4 semanas & & & \\
Prejuizo na participacao escolar & & & \\
Inicio na infancia (antes dos 6 anos tipico) & & & \\
\bottomrule\end{tabular}
\subsection*{Como investigar (tarefas sistematicas)}
\begin{itemize}
\item Registro de sintomas em dias de entrada
\item Entrevista com pais (contexto e gatilhos)
\item Observacao da adaptacao progressiva
\end{itemize}
\subsection*{Exclusoes e confundidores}
\begin{itemize}
\item Falta de rotina de despedida
\item Evento estressor recente (luto, mudanca)
\item Doenca fisica
\end{itemize}
\subsection*{Instrumentos validados para confirmacao (referencia -- aplicar somente por psicologo, conforme SATEPSI; nao reproduzidos aqui)}
\begin{itemize}
\item Escalas de ansiedade infantil (SCARED/CAS -- referencia)
\item Entrevista clinica (psicologo)
\end{itemize}
\subsection*{Encaminhamento}
\begin{itemize}
\item Psicologo (terapia cognitivo-comportamental)
\item Fonoaudiologo nao indicado (sem linguagem)
\end{itemize}
\subsection*{Adaptacoes pedagogicas iniciais (sem rotulo)}
\begin{itemize}
\item Ritual de despedida curto e consistente
\item Objeto de conforto (foto, bichinho)
\item Chegada gradual com profissional de apoio
\item Reforco positivo apos periodos de permanencia
\item Comunicacao com a familia sobre avancos
\end{itemize}
\subsection*{Bibliografia basica}
\begin{itemize}
\item DSM-5-TR (2023). Artmed.
\item Caballo (2007), Manual de TEAs e de transtornos fobicos. VCH.
\end{itemize}
\avisolegal


\chapter{Ansiedade Generalizada / de Desempenho}
\section*{Dominio: Socioemocional}
\section*{Ficha de rastreio sistematico -- confirmacao}
\subsection*{Sinais de alerta (marque: S = sim, AV = as vezes, N = nao)}
\begin{tabular}{@{}p{11.5cm}ccc@{}}\toprule
Sinal observado & S & AV & N \\\midrule
Preocupacoes excessivas sobre desempenho escolar & & & \\
Dificuldade de controlar a preocupacao & & & \\
Sintomas fisicos (tensao, cansaco, insonia) sem causa medica & & & \\
Busca constante de aprovacao & & & \\
Evita tarefas com medo de errar & & & \\
Perguntas repetitivas sobre o futuro & & & \\
Irritabilidade sob cobranca & & & \\
Dificuldade de relaxar & & & \\
Sinais ha pelo menos 6 meses & & & \\
Prejuizo no rendimento e no bem-estar & & & \\
\bottomrule\end{tabular}
\subsection*{Como investigar (tarefas sistematicas)}
\begin{itemize}
\item Entrevista com a crianca (percepcoes)
\item Escala de ansiedade (SCARED -- referencia)
\item Observacao em situacoes de avaliacao
\end{itemize}
\subsection*{Exclusoes e confundidores}
\begin{itemize}
\item Cobranca familiar intensa
\item TDAH (procrastinacao por attencao)
\item Perfeccionismo adaptativo sem prejuizo
\end{itemize}
\subsection*{Instrumentos validados para confirmacao (referencia -- aplicar somente por psicologo, conforme SATEPSI; nao reproduzidos aqui)}
\begin{itemize}
\item SCARED/CAS (uso profissional)
\item Avaliacao clinica (psicologo)
\end{itemize}
\subsection*{Encaminhamento}
\begin{itemize}
\item Psicologo (TCC) se prejuizo funcional
\item Pediatra para sintomas somaticos
\end{itemize}
\subsection*{Adaptacoes pedagogicas iniciais (sem rotulo)}
\begin{itemize}
\item Ambiente de erro permitido e seguro (erro como aprendizagem)
\item Instrucoes claras com objetivos pequenos
\item Tecnicas de respiracao antes de provas
\item Reforco de esforco e processo, nao so do resultado
\item Reduzir pressao de tempo em avaliacoes quando indicado
\end{itemize}
\subsection*{Bibliografia basica}
\begin{itemize}
\item DSM-5-TR (2023). Artmed.
\item Stallard (2010), Ansiedade infantil. Artmed.
\end{itemize}
\avisolegal


\chapter{Fobia Escolar / Ansiedade Social}
\section*{Dominio: Socioemocional}
\section*{Ficha de rastreio sistematico -- confirmacao}
\subsection*{Sinais de alerta (marque: S = sim, AV = as vezes, N = nao)}
\begin{tabular}{@{}p{11.5cm}ccc@{}}\toprule
Sinal observado & S & AV & N \\\midrule
Recusa persistente de ir a escola com reacao intensa & & & \\
Angustia antecipatoria ja na noite anterior & & & \\
Sintomas somaticos fortes na entrada (vomito, dor) & & & \\
Medo intenso de situacoes de exposicao (falar na frente) & & & \\
Evita interacoes sociais com medo de avaliacao & & & \\
Fica mudo/a ou colapsa em chamadas orais & & & \\
Dificuldade de olhar/responder a adultos & & & \\
Ausencias frequentes por 'doenca' sem causa & & & \\
Escolarizacao em risco & & & \\
Sintomas presentes em casa com familiares conhecidos (contraste) & & & \\
\bottomrule\end{tabular}
\subsection*{Como investigar (tarefas sistematicas)}
\begin{itemize}
\item Registro de padrao de ausencias e sintomas
\item Entrevista com pais e crianca
\item Observacao em situacoes de exposicao
\end{itemize}
\subsection*{Exclusoes e confundidores}
\begin{itemize}
\item Bullying (investigar antes)
\item Ansiedade de separacao
\item Dificuldades academicas como gatilho
\end{itemize}
\subsection*{Instrumentos validados para confirmacao (referencia -- aplicar somente por psicologo, conforme SATEPSI; nao reproduzidos aqui)}
\begin{itemize}
\item Escalas de ansiedade social
\item Avaliacao clinica (psicologo)
\end{itemize}
\subsection*{Encaminhamento}
\begin{itemize}
\item Psicologo (TCC com exposicao gradual)
\item Contato com rede de saude se absenteismo escolar
\end{itemize}
\subsection*{Adaptacoes pedagogicas iniciais (sem rotulo)}
\begin{itemize}
\item Reentrada gradual com plano de exposicao
\item Tarefa de exposicao oral pequena e crescente (dupla, trio, grupo)
\item Apoio de profissional em momentos-chave
\item Reforco positivo de frequencia
\item Envolver a familia em plano conjunto
\end{itemize}
\subsection*{Bibliografia basica}
\begin{itemize}
\item DSM-5-TR (2023). Artmed.
\item Heyne \& Rollings (2003). School refusal. Wiley.
\end{itemize}
\avisolegal


\chapter{Depressao Infantil}
\section*{Dominio: Socioemocional}
\section*{Ficha de rastreio sistematico -- confirmacao}
\subsection*{Sinais de alerta (marque: S = sim, AV = as vezes, N = nao)}
\begin{tabular}{@{}p{11.5cm}ccc@{}}\toprule
Sinal observado & S & AV & N \\\midrule
Tristeza ou irritabilidade persistente & & & \\
Perda de interesse em atividades antes prazerosas & & & \\
Queda no rendimento e na concentracao & & & \\
Fadiga, hipoatividade ou agitacao & & & \\
Mudanca de apetite/sono & & & \\
Autodesvalorizacao ('sou incapaz') & & & \\
Isolamento social & & & \\
Lentidao motora ou fala reduzida & & & \\
Sintomas ha pelo menos 2 semanas & & & \\
Historia de perdas/eventos estressores & & & \\
\bottomrule\end{tabular}
\subsection*{Como investigar (tarefas sistematicas)}
\begin{itemize}
\item Entrevista com a crianca e os pais
\item Escala de depressao infantil (CDI -- referencia)
\item Observacao na rotina escolar
\end{itemize}
\subsection*{Exclusoes e confundidores}
\begin{itemize}
\item Luto reativo recente
\item Ansiedade com apatia secundaria
\item Hipotireoidismo/anemia (avaliar medico)
\end{itemize}
\subsection*{Instrumentos validados para confirmacao (referencia -- aplicar somente por psicologo, conforme SATEPSI; nao reproduzidos aqui)}
\begin{itemize}
\item CDI (uso profissional)
\item Avaliacao clinica/psiquiatrica
\end{itemize}
\subsection*{Encaminhamento}
\begin{itemize}
\item Psicologo/psiquiatra infantil (urgente se risco)
\item Acolhimento e rede de apoio
\end{itemize}
\subsection*{Adaptacoes pedagogicas iniciais (sem rotulo)}
\begin{itemize}
\item Acolhimento individual breve diario
\item Tarefas pequenas com alta chance de sucesso
\item Reconhecer esforco e presenca
\item Envolver a familia e a rede (ESF/CRAS)
\item Proibido rotular; encaminhamento sigiloso
\end{itemize}
\subsection*{Bibliografia basica}
\begin{itemize}
\item DSM-5-TR (2023). Artmed.
\item Bahls (2002), Depressao na infancia e adolescencia. UFPR.
\end{itemize}
\avisolegal


\chapter{TOC Infantil}
\section*{Dominio: Socioemocional}
\section*{Ficha de rastreio sistematico -- confirmacao}
\subsection*{Sinais de alerta (marque: S = sim, AV = as vezes, N = nao)}
\begin{tabular}{@{}p{11.5cm}ccc@{}}\toprule
Sinal observado & S & AV & N \\\midrule
Preocupacoes repetitivas que a crianca tenta ignorar & & & \\
Rituais de limpeza, ordem, verificacao ou contagem & & & \\
Necessidade de simetria nas atividades & & & \\
Perda de tempo significativa com ritual & & & \\
Angustia se impedido de realizar ritual & & & \\
Perguntas repetitivas ('e se...?') & & & \\
Sinais comecam na infancia (media 7-12 anos) & & & \\
Prejuizo academico e social & & & \\
Tentativas de esconder rituais & & & \\
Sintomas consomem mais de 1h/dia & & & \\
\bottomrule\end{tabular}
\subsection*{Como investigar (tarefas sistematicas)}
\begin{itemize}
\item Entrevista clinica cuidadosa
\item Observacao de rituais na sala
\item Excluir transtornos do desenvolvimento
\end{itemize}
\subsection*{Exclusoes e confundidores}
\begin{itemize}
\item Ansiedade geral com verificacao leve
\item Tiques simples
\item Rotinas normais de infancia (transitorias)
\end{itemize}
\subsection*{Instrumentos validados para confirmacao (referencia -- aplicar somente por psicologo, conforme SATEPSI; nao reproduzidos aqui)}
\begin{itemize}
\item Avaliacao clinica/psiquiatrica
\item Escalas especificas (CY-BOCS -- referencia)
\end{itemize}
\subsection*{Encaminhamento}
\begin{itemize}
\item Psiquiatra/psicologo (TCC com prevencao de resposta)
\end{itemize}
\subsection*{Adaptacoes pedagogicas iniciais (sem rotulo)}
\begin{itemize}
\item Flexibilizar prazos e exigencia de limpeza/ordem
\item Sem punicao por rituais; registrar o impacto
\item Plano de resposta a estereotipias com profissional
\item Ambiente estruturado com atividades que engajam a atencao
\end{itemize}
\subsection*{Bibliografia basica}
\begin{itemize}
\item DSM-5-TR (2023). Artmed.
\item March \& Mulle (1998). OCD in children. Guilford.
\end{itemize}
\avisolegal


\chapter{Mutismo Seletivo}
\section*{Dominio: Linguagem/Socioemocional}
\section*{Ficha de rastreio sistematico -- confirmacao}
\subsection*{Sinais de alerta (marque: S = sim, AV = as vezes, N = nao)}
\begin{tabular}{@{}p{11.5cm}ccc@{}}\toprule
Sinal observado & S & AV & N \\\midrule
Fala em casa e nao fala na escola & & & \\
Ausencia de fala em situacoes especificas por mais de 1 mes & & & \\
Nao fala mesmo quando solicitado ou sob pressao & & & \\
Usa gestos, acenos e escrita para responder & & & \\
Contato visual e postura podem ser normais ou inibidos & & & \\
Aparece ao entrar na escola (transicao) & & & \\
Sinais de ansiedade social associados & & & \\
Fala com selecao de pessoas e locais & & & \\
Crianca e comunicativa em ambiente seguro & & & \\
Nao explicado por TEA ou TDL & & & \\
\bottomrule\end{tabular}
\subsection*{Como investigar (tarefas sistematicas)}
\begin{itemize}
\item Entrevista com pais (contextos em que fala)
\item Observacao em diferentes situacoes (pequeno grupo, recreio)
\item Excluir TDL e TEA (comunicacao global)
\end{itemize}
\subsection*{Exclusoes e confundidores}
\begin{itemize}
\item Timidez normal de adaptacao (menos de 1 mes)
\item TDL/TEA
\item Trauma/refugio
\end{itemize}
\subsection*{Instrumentos validados para confirmacao (referencia -- aplicar somente por psicologo, conforme SATEPSI; nao reproduzidos aqui)}
\begin{itemize}
\item Avaliacao clinica (psicologo)
\item Fonoaudiologo para excluir TDL
\end{itemize}
\subsection*{Encaminhamento}
\begin{itemize}
\item Psicologo (terapia com exposicao gradual)
\item Fonoaudiologo se comorbidade com TDL
\end{itemize}
\subsection*{Adaptacoes pedagogicas iniciais (sem rotulo)}
\begin{itemize}
\item Nao pressionar para falar; aceitar comunicacao alternativa
\item Situacoes de fala em dupla com colega seguro
\item Reforco de qualquer forma de comunicacao
\item Atividades de resposta nao-verbal inicial (apontar, cartoes)
\item Combinado com a familia sobre a escola
\end{itemize}
\subsection*{Bibliografia basica}
\begin{itemize}
\item DSM-5-TR (2023). Artmed.
\item Kearney (2010). Helping children with selective mutism. Oxford.
\end{itemize}
\avisolegal


\chapter{TPAC (processamento auditivo)}
\section*{Dominio: Processamento Auditivo}
\section*{Ficha de rastreio sistematico -- confirmacao}
\subsection*{Sinais de alerta (marque: S = sim, AV = as vezes, N = nao)}
\begin{tabular}{@{}p{11.5cm}ccc@{}}\toprule
Sinal observado & S & AV & N \\\midrule
Pedidos frequentes de repeticao & & & \\
Dificuldade em ambiente ruidoso & & & \\
Confusoes de palavras parecidas ao ouvir & & & \\
Seguir instrucoes longas com erros & & & \\
Dificuldade de aprender rimas e musicas & & & \\
Leitura e ortografia prejudicadas por limite fonologico & & & \\
Atraso em compreender piadas/ironias orais & & & \\
Audicao periferica normal (avaliacao otorrinolaringologica) & & & \\
Melhora com apoio visual & & & \\
Cansa ao ouvir explicacoes longas & & & \\
\bottomrule\end{tabular}
\subsection*{Como investigar (tarefas sistematicas)}
\begin{itemize}
\item Triagem auditiva formal (OTORRINO) para excluir perda
\item Provas de memoria sequencial auditiva
\item Observar em sala com e sem ruido
\end{itemize}
\subsection*{Exclusoes e confundidores}
\begin{itemize}
\item Perda auditiva/TDAH (desatencao auditiva)
\item TDL
\item Ambiente escolar ruidoso
\end{itemize}
\subsection*{Instrumentos validados para confirmacao (referencia -- aplicar somente por psicologo, conforme SATEPSI; nao reproduzidos aqui)}
\begin{itemize}
\item Avaliacao fonoaudiologica de PA (somente apos audicao normal)
\item Testes de processamento auditivo (fonoaudiologo)
\end{itemize}
\subsection*{Encaminhamento}
\begin{itemize}
\item Fonoaudiologo especialista em PA
\item OTORRINO primeiro
\end{itemize}
\subsection*{Adaptacoes pedagogicas iniciais (sem rotulo)}
\begin{itemize}
\item Assento proximo ao professor
\item Instrucoes curtas e checagem da compreensao
\item Reduzir ruido de fundo na sala
\item Apoio visual em todas as instrucoes orais
\item Repetir e reescrever instrucoes de provas
\end{itemize}
\subsection*{Bibliografia basica}
\begin{itemize}
\item Ferreira (2004), Tratado de fonoaudiologia. Roca.
\item Pereira \& Schochat (2011). PA. Pulso.
\end{itemize}
\avisolegal


\chapter{TDL (Transtorno do Desenvolvimento da Linguagem)}
\section*{Dominio: Linguagem Oral}
\section*{Ficha de rastreio sistematico -- confirmacao}
\subsection*{Sinais de alerta (marque: S = sim, AV = as vezes, N = nao)}
\begin{tabular}{@{}p{11.5cm}ccc@{}}\toprule
Sinal observado & S & AV & N \\\midrule
Atraso no vocabulario e nas combinacoes de frases & & & \\
Frases curtas ou agramaticais para a idade & & & \\
Dificuldade de narrar e relatar fatos & & & \\
Problemas de compreensao de ordens complexas & & & \\
Uso pobre de conectivos (porque, mas) & & & \\
Dificuldade de aprender novas palavras & & & \\
Prejuizo na leitura/escrita decorrente & & & \\
Sem perda auditiva, TEA ou DI que expliquem & & & \\
Melhora com suporte visual & & & \\
Historico de fala tardia & & & \\
\bottomrule\end{tabular}
\subsection*{Como investigar (tarefas sistematicas)}
\begin{itemize}
\item Avaliacao de linguagem (fonoaudiologo)
\item Prova de vocabulario e narrativa
\item Observacao do discurso em sala
\end{itemize}
\subsection*{Exclusoes e confundidores}
\begin{itemize}
\item Perda auditiva
\item TEA (pragmatica comprometida global)
\item DI (prejuizo global)
\item Mutismo seletivo
\end{itemize}
\subsection*{Instrumentos validados para confirmacao (referencia -- aplicar somente por psicologo, conforme SATEPSI; nao reproduzidos aqui)}
\begin{itemize}
\item ABFW (referencia; fonoaudiologo)
\item Avaliacao de linguagem formal
\end{itemize}
\subsection*{Encaminhamento}
\begin{itemize}
\item Fonoaudiologo (terapia de linguagem intensiva)
\end{itemize}
\subsection*{Adaptacoes pedagogicas iniciais (sem rotulo)}
\begin{itemize}
\item Linguagem simplificada e concreta com apoio visual
\item Modelar frases completas sem corrigir publicamente
\item Antecipar vocabulario novo
\item Verificar compreensao com perguntas especificas
\item Pares mais fluentes em trabalhos em grupo
\end{itemize}
\subsection*{Bibliografia basica}
\begin{itemize}
\item Bishop et al. (2017), TDL. JCPP.
\item Hage \& Guimaraes (2006). Disturbios de linguagem. Pulso.
\end{itemize}
\avisolegal


\chapter{Desvio Fonologico (dislalia)}
\section*{Dominio: Linguagem Oral}
\section*{Ficha de rastreio sistematico -- confirmacao}
\subsection*{Sinais de alerta (marque: S = sim, AV = as vezes, N = nao)}
\begin{tabular}{@{}p{11.5cm}ccc@{}}\toprule
Sinal observado & S & AV & N \\\midrule
Trocas de fonemas na fala (tato/gato) & & & \\
Omissao de silabas ou fonemas & & & \\
Fala inteligivel apenas para a familia & & & \\
Dificuldade de articulacao de vibrante (r) & & & \\
Erros fonologicos alem da idade esperada (6 anos: eliminar maioria) & & & \\
Afeta escrita (troca correspondente) & & & \\
Vergonha de falar em publico & & & \\
Sem perda auditiva ou alteracao anatomica & & & \\
Resposta inconsistente (mesma palavra, erros diferentes) & & & \\
Persistencia apos ensino explicito & & & \\
\bottomrule\end{tabular}
\subsection*{Como investigar (tarefas sistematicas)}
\begin{itemize}
\item Avaliacao fonoaudiologica articulatoria
\item Prova de nomeacao de figuras
\item Excluir perda auditiva e PA
\end{itemize}
\subsection*{Exclusoes e confundidores}
\begin{itemize}
\item Perda auditiva
\item TDL mais amplo
\item Variacao regional da fala (nao patologia)
\end{itemize}
\subsection*{Instrumentos validados para confirmacao (referencia -- aplicar somente por psicologo, conforme SATEPSI; nao reproduzidos aqui)}
\begin{itemize}
\item ABFW (prova fonologica; fonoaudiologo)
\item Avaliacao de PA se necessario
\end{itemize}
\subsection*{Encaminhamento}
\begin{itemize}
\item Fonoaudiologo (terapia fonologica)
\end{itemize}
\subsection*{Adaptacoes pedagogicas iniciais (sem rotulo)}
\begin{itemize}
\item Modelar corretamente sem pedir repeticao publica
\item Atividades de consciencia fonologica ludicas
\item Nao ridicularizar; reforcar tentativas
\item Registrar erros para a fonoaudiologa
\item Leitura de palavras com o fonema-alvo
\end{itemize}
\subsection*{Bibliografia basica}
\begin{itemize}
\item Wertzner (2000), Fonologia. Pulso.
\item Lamprecht (2004), Aquisicao fonologica. Artmed.
\end{itemize}
\avisolegal


\chapter{Gagueira (tartamudez)}
\section*{Dominio: Fala}
\section*{Ficha de rastreio sistematico -- confirmacao}
\subsection*{Sinais de alerta (marque: S = sim, AV = as vezes, N = nao)}
\begin{tabular}{@{}p{11.5cm}ccc@{}}\toprule
Sinal observado & S & AV & N \\\midrule
Repeticoes de silabas/sons (mu-mu-muito) & & & \\
Prolongamentos (m====aior) & & & \\
Bloqueios na fala (fica travado) & & & \\
Esforco fisico visivel ao falar & & & \\
Evita palavras ou situacoes de fala & & & \\
Tensao facial ou piscadas frequentes & & & \\
Gagueira com duracao superior a 6 meses & & & \\
Historia familiar de gagueira & & & \\
Sinais de frustracao/vergonha & & & \\
Fluencia varia conforme situacao (leitura, pressa) & & & \\
\bottomrule\end{tabular}
\subsection*{Como investigar (tarefas sistematicas)}
\begin{itemize}
\item Observacao em diferentes situacoes de fala
\item Registro de tipologia (repeticoes, prolongamentos, bloqueios)
\item Entrevista com a familia
\end{itemize}
\subsection*{Exclusoes e confundidores}
\begin{itemize}
\item Disfluencias normais do desenvolvimento (fase prescolare ate 6 anos, sem tensao)
\item Fala rapida por ansiedade
\end{itemize}
\subsection*{Instrumentos validados para confirmacao (referencia -- aplicar somente por psicologo, conforme SATEPSI; nao reproduzidos aqui)}
\begin{itemize}
\item Avaliacao fonoaudiologica de fluencia
\item Encaminhamento precoce (quanto antes, melhor)
\end{itemize}
\subsection*{Encaminhamento}
\begin{itemize}
\item Fonoaudiologo especialista em fluencia (preferencialmente em ate 12 meses do inicio)
\end{itemize}
\subsection*{Adaptacoes pedagogicas iniciais (sem rotulo)}
\begin{itemize}
\item Nao interromper, corrigir ou completar a fala
\item Dar tempo (esperar sem pressa)
\item Reduzir pressao de velocidade em leitura oral
\item Falar com a crianca em ritmo pausado como modelo
\item Evitar humilhacoes publicas; reforcar conteudo
\end{itemize}
\subsection*{Bibliografia basica}
\begin{itemize}
\item Andrade (2006), Fluencia. Roca.
\item Nippold (2011), Stuttering. Psychology Press.
\end{itemize}
\avisolegal


\chapter{Apraxia de Fala Infantil}
\section*{Dominio: Fala}
\section*{Ficha de rastreio sistematico -- confirmacao}
\subsection*{Sinais de alerta (marque: S = sim, AV = as vezes, N = nao)}
\begin{tabular}{@{}p{11.5cm}ccc@{}}\toprule
Sinal observado & S & AV & N \\\midrule
Fala muito limitada ou ininteligivel apesar de esforco & & & \\
Erros inconsistentes e variaveis na mesma palavra & & & \\
Dificuldade de organizar sequencias de movimentos da fala & & & \\
Piora progressiva em frases longas & & & \\
Atraso no desenvolvimento motor oral (soprar, estalar labios) & & & \\
Historia de dificuldades na alimentacao & & & \\
Compreensao melhor que producao & & & \\
Melhora com apoio visual/silabas modeladas lentamente & & & \\
Sem alteracao neurologica evidente (excluir paralisia) & & & \\
Frustracao comunicativa intensa & & & \\
\bottomrule\end{tabular}
\subsection*{Como investigar (tarefas sistematicas)}
\begin{itemize}
\item Avaliacao fonoaudiologica detalhada
\item Prova de diadococinesia (repeticao rapida)
\item Gravacao de amostra de fala
\end{itemize}
\subsection*{Exclusoes e confundidores}
\begin{itemize}
\item Desvio fonologico (erros consistentes)
\item TDL
\item Dispraxia geral/TDC
\end{itemize}
\subsection*{Instrumentos validados para confirmacao (referencia -- aplicar somente por psicologo, conforme SATEPSI; nao reproduzidos aqui)}
\begin{itemize}
\item Avaliacao fonoaudiologica (especialista em motor da fala)
\item Avaliacao neurologica se necessario
\end{itemize}
\subsection*{Encaminhamento}
\begin{itemize}
\item Fonoaudiologo (terapia motora intensiva)
\end{itemize}
\subsection*{Adaptacoes pedagogicas iniciais (sem rotulo)}
\begin{itemize}
\item Comunicacao alternativa temporaria se frustracao alta
\item Palavras-alvo curtas modeladas lentamente
\item Elogiar tentativas, nao exigir perfeicao
\item Antecipar e verbalizar pelo aluno em situacoes de estresse
\end{itemize}
\subsection*{Bibliografia basica}
\begin{itemize}
\item Strand et al. (2013). CAS. ASHA.
\item Dodd (2005), Differential diagnosis. Whurr.
\end{itemize}
\avisolegal


\chapter{TDC / Dispraxia (coordenacao)}
\section*{Dominio: Motricidade}
\section*{Ficha de rastreio sistematico -- confirmacao}
\subsection*{Sinais de alerta (marque: S = sim, AV = as vezes, N = nao)}
\begin{tabular}{@{}p{11.5cm}ccc@{}}\toprule
Sinal observado & S & AV & N \\\midrule
Clumsiness: esbarra, derruba objetos com frequencia & & & \\
Dificuldade de aprender habilidades motoras (pular corda, andar de bicicleta) & & & \\
Escrita muito lenta e desajeitada (ver ficha 3) & & & \\
Dificuldade de abotoar, amarrar, usar talheres & & & \\
Quedas frequentes sem causa & & & \\
Desempenho motor muito abaixo do esperado para a idade & & & \\
Dificuldade de imitar movimentos & & & \\
Evita atividades fisicas por constrangimento & & & \\
Interfere em atividades da vida diaria e escolar & & & \\
Sinais desde cedo (nao adquiridos por falta de pratica) & & & \\
\bottomrule\end{tabular}
\subsection*{Como investigar (tarefas sistematicas)}
\begin{itemize}
\item Avaliacao psicomotora (MABC-2 -- referencia)
\item Testes de destreza manual e equilibrio
\item Observacao em educacao fisica e recreio
\end{itemize}
\subsection*{Exclusoes e confundidores}
\begin{itemize}
\item Falta de pratica/escolarizacao
\item Problema visual
\item Doenca neurologica (excluir)
\end{itemize}
\subsection*{Instrumentos validados para confirmacao (referencia -- aplicar somente por psicologo, conforme SATEPSI; nao reproduzidos aqui)}
\begin{itemize}
\item MABC-2 (uso profissional)
\item Avaliacao psicomotora/terapeuta ocupacional
\end{itemize}
\subsection*{Encaminhamento}
\begin{itemize}
\item Terapia ocupacional/psicomotricidade
\item Atividade fisica adaptada
\end{itemize}
\subsection*{Adaptacoes pedagogicas iniciais (sem rotulo)}
\begin{itemize}
\item Adaptar instrumentos (lapis grosso, antiderrapante)
\item Dividir tarefas motoras em etapas pequenas
\item Tempo extra em copias
\item Reforcar desempenho pelo esforco
\item Atividades fisicas alternativas de sucesso
\end{itemize}
\subsection*{Bibliografia basica}
\begin{itemize}
\item DSM-5-TR (2023). Artmed.
\item Fonseca (2012), Manual de avaliacao psicomotora. Ancora.
\end{itemize}
\avisolegal


\chapter{Tiques e Tourette}
\section*{Dominio: Comportamento/Motor}
\section*{Ficha de rastreio sistematico -- confirmacao}
\subsection*{Sinais de alerta (marque: S = sim, AV = as vezes, N = nao)}
\begin{tabular}{@{}p{11.5cm}ccc@{}}\toprule
Sinal observado & S & AV & N \\\midrule
Movimentos rapidos involuntarios (piscar, sacudir cabeca) & & & \\
Vocalizacoes involuntarias (fungar, pigarrear, palavras) & & & \\
Tiques mudam ao longo do tempo & & & \\
Piora com estresse, fadiga e ansiedade & & & \\
Crianca consegue suprimir temporariamente & & & \\
Inicio tipico entre 5-10 anos & & & \\
Sinais ha mais de 1 ano para Tourette (motor + vocal) & & & \\
Sensacao de alivio apos tique & & & \\
Prejuizo social e academico & & & \\
Tiques nao explicados por medicacao ou substancia & & & \\
\bottomrule\end{tabular}
\subsection*{Como investigar (tarefas sistematicas)}
\begin{itemize}
\item Observacao sistematica em sala
\item Entrevista com pais (duracao, tipos)
\item Registro em diferentes contextos
\end{itemize}
\subsection*{Exclusoes e confundidores}
\begin{itemize}
\item Habitos motores normais
\item TDAH com agitacao
\item TOC (rituais vs tiques)
\end{itemize}
\subsection*{Instrumentos validados para confirmacao (referencia -- aplicar somente por psicologo, conforme SATEPSI; nao reproduzidos aqui)}
\begin{itemize}
\item Avaliacao clinica/neurologica/psiquiatrica
\item Escalas de tiques (YGTSS -- referencia)
\end{itemize}
\subsection*{Encaminhamento}
\begin{itemize}
\item Neurologista/psiquiatra infantil se prejuizo
\item Psicologo (TCC-HRT) se indicado
\end{itemize}
\subsection*{Adaptacoes pedagogicas iniciais (sem rotulo)}
\begin{itemize}
\item Nao chamar atencao para o tique; nao punir
\item Reduzir estresse e pressao (gatilhos)
\item Permitir pausas de saida quando necessario
\item Informar colegas de forma natural se a crianca concordar
\item Coordenacao com medico
\end{itemize}
\subsection*{Bibliografia basica}
\begin{itemize}
\item DSM-5-TR (2023). Artmed.
\item Leckman et al. (2010). Tourette. Guilford.
\end{itemize}
\avisolegal


\chapter{Ausencias (sinais para encaminhamento medico)}
\section*{Dominio: Saude/Neuro}
\section*{Ficha de rastreio sistematico -- confirmacao}
\subsection*{Sinais de alerta (marque: S = sim, AV = as vezes, N = nao)}
\begin{tabular}{@{}p{11.5cm}ccc@{}}\toprule
Sinal observado & S & AV & N \\\midrule
Piscadelas ou olhar fixo sem resposta por segundos & & & \\
Paradas breves da atividade com retomada & & & \\
Automatsmos (movimentos de mastigacao) durante episodio & & & \\
Quedas de cabeca sem causa aparente & & & \\
Perda de urina em horarios incomuns & & & \\
Nao responde a chamados durante o episodio & & & \\
Lapsos repetidos de 'sonho acordado' & & & \\
Parentes descrevem 'desligamentos' & & & \\
Ausencias podem ser confundidas com desatencao & & & \\
Prejuizo de aprendizagem por episodios nao percebidos & & & \\
\bottomrule\end{tabular}
\subsection*{Como investigar (tarefas sistematicas)}
\begin{itemize}
\item Registrar com data/hora/doracao dos episodios
\item Filmar com celular (com autorizacao) para o medico
\item Observar resposta a estimulos durante o episodio
\end{itemize}
\subsection*{Exclusoes e confundidores}
\begin{itemize}
\item Desatencao (TDAH) -- o aluno RESPONDE ao chamado
\item Sonolencia por sono ruim
\item Transtorno de ausencia precisa confirmacao medica
\end{itemize}
\subsection*{Instrumentos validados para confirmacao (referencia -- aplicar somente por psicologo, conforme SATEPSI; nao reproduzidos aqui)}
\begin{itemize}
\item Encaminhamento ao pediatra/neurologista (NUNCA tratar na escola)
\item EEG se indicado pelo medico
\end{itemize}
\subsection*{Encaminhamento}
\begin{itemize}
\item Pediatra/neurologista pediatrico (URGENTE se episodios frequentes)
\end{itemize}
\subsection*{Adaptacoes pedagogicas iniciais (sem rotulo)}
\begin{itemize}
\item NAO dar diagnostico escolar
\item Registrar episodios e repassar a familia/medico
\item Avaliar risco de queda
\item Nao deixar a crianca em altura ou atividades perigosas ate avaliacao
\item Manter rotina escolar com apoio
\end{itemize}
\subsection*{Bibliografia basica}
\begin{itemize}
\item DSM-5-TR (2023). Artmed.
\item Camfield \& Camfield (2005). Pediatric epilepsy. INPE.
\end{itemize}
\avisolegal


\chapter{Deficiencia Intelectual leve (sinais escolares)}
\section*{Dominio: Desenvolvimento Global}
\section*{Ficha de rastreio sistematico -- confirmacao}
\subsection*{Sinais de alerta (marque: S = sim, AV = as vezes, N = nao)}
\begin{tabular}{@{}p{11.5cm}ccc@{}}\toprule
Sinal observado & S & AV & N \\\midrule
Atraso global de aprendizagem em todas as areas & & & \\
Dificuldade de generalizar aprendizagens (contexto a contexto) & & & \\
Linguagem simples, concreta e literal & & & \\
Dificuldade de resolver problemas novos & & & \\
Memoria de trabalho reduzida & & & \\
Dependencia de apoio para tarefas cotidianas & & & \\
Dificuldade de raciocinio logico e abstrato & & & \\
Adaptacao escolar necessaria em ritmo proprio & & & \\
Sem explicacao por privacao sensorial grave & & & \\
Sinais desde a infancia & & & \\
\bottomrule\end{tabular}
\subsection*{Como investigar (tarefas sistematicas)}
\begin{itemize}
\item Avaliacao interdisciplinar (pedagogica + psicologica)
\item Observacao de autonomia e aprendizagem
\item Historia do desenvolvimento
\end{itemize}
\subsection*{Exclusoes e confundidores}
\begin{itemize}
\item Falta de escolarizacao/bilinguismo
\item TDAH (atencao afeta todas as areas)
\item TDL
\item Privacao sociocultural (pobreza, negligencia)
\end{itemize}
\subsection*{Instrumentos validados para confirmacao (referencia -- aplicar somente por psicologo, conforme SATEPSI; nao reproduzidos aqui)}
\begin{itemize}
\item WISC-V / avaliacao intelectual (psicologo, SATEPSI)
\item Escalas adaptativas (Vineland -- referencia)
\end{itemize}
\subsection*{Encaminhamento}
\begin{itemize}
\item Psicologo (avaliacao formal) e medico para causas
\item AEE (Atendimento Educacional Especializado)
\end{itemize}
\subsection*{Adaptacoes pedagogicas iniciais (sem rotulo)}
\begin{itemize}
\item Curriculo funcional com metas pequenas e concretas
\item Ensino com muitos exemplos e pratica
\item Avaliacao oral/adaptada
\item Apoio de AEE e plano individualizado
\item Instrumentos concretos e material visual
\end{itemize}
\subsection*{Bibliografia basica}
\begin{itemize}
\item DSM-5-TR (2023). Artmed.
\item Mendonca (2012). DI e escola. Juruá.
\end{itemize}
\avisolegal


\chapter{Altas Habilidades / Superdotacao}
\section*{Dominio: Desenvolvimento/Curriculo}
\section*{Ficha de rastreio sistematico -- confirmacao}
\subsection*{Sinais de alerta (marque: S = sim, AV = as vezes, N = nao)}
\begin{tabular}{@{}p{11.5cm}ccc@{}}\toprule
Sinal observado & S & AV & N \\\midrule
Aprendizagem muito rapida e facil em uma ou mais areas & & & \\
Vocabulario avancado e curiosidade intensa & & & \\
Memoria e raciocinio acima da media & & & \\
Interesse intenso e aprofundado por um tema & & & \\
Faz perguntas desafiadoras & & & \\
Perfeccionismo que gera frustracao & & & \\
Pode apresentar tedio e desatencao em tarefas repetitivas & & & \\
Sensibilidade emocional aumentada & & & \\
Precisa de desafios para engajar & & & \\
Criatividade e solucao original de problemas & & & \\
\bottomrule\end{tabular}
\subsection*{Como investigar (tarefas sistematicas)}
\begin{itemize}
\item Observacao de desempenho e interesse
\item Anamnese com familia (marcos precoces)
\item Avaliacao psicologica se indicado
\end{itemize}
\subsection*{Exclusoes e confundidores}
\begin{itemize}
\item TDAH (desatento por tedio vs desatento por deficit)
\item Ansiedade por pressao
\item Comportamento de aluno 'esperto' sem excepcionalidade
\end{itemize}
\subsection*{Instrumentos validados para confirmacao (referencia -- aplicar somente por psicologo, conforme SATEPSI; nao reproduzidos aqui)}
\begin{itemize}
\item Avaliacao psicologica (psicologo, SATEPSI)
\item Testes de inteligencia se indicado (WISC-V)
\end{itemize}
\subsection*{Encaminhamento}
\begin{itemize}
\item Psicologo (avaliacao formal) se necessario
\item Programas de enriquecimento/NAAH/S
\end{itemize}
\subsection*{Adaptacoes pedagogicas iniciais (sem rotulo)}
\begin{itemize}
\item Enriquecimento curricular (projetos, mentorias)
\item Tarefas de nivel avancado (nao so mais do mesmo)
\item Grupo de pares com interesses comuns
\item Apoio socioemocional para perfeccionismo
\item Plano de trabalho individualizado
\end{itemize}
\subsection*{Bibliografia basica}
\begin{itemize}
\item Renzulli (2005), Three-ring conception. Kappan.
\item Alencar \& Fleith (2001). Superdotados. EPU.
\end{itemize}
\avisolegal


\chapter{Perfil de Aprendizagem Nao-Verbal (TANV)}
\section*{Dominio: Visomotor/Pragmatica}
\section*{Ficha de rastreio sistematico -- confirmacao}
\subsection*{Sinais de alerta (marque: S = sim, AV = as vezes, N = nao)}
\begin{tabular}{@{}p{11.5cm}ccc@{}}\toprule
Sinal observado & S & AV & N \\\midrule
TDE verbal/leitura boa, matematica visual-praxia pobre & & & \\
Dificuldade de interpretar expressoes e linguagem corporal & & & \\
Problemas de coordenacao fina (ver ficha 22) & & & \\
Dificuldade de organizacao visuoespacial (perder-se, copia ruim) & & & \\
Pragmatica social desajeitada & & & \\
Compreensao literal elevada, ironia dificil & & & \\
Dificuldade com conceitos abstratos em matematica (valor posicional) & & & \\
Rotina inflexivel com ansiedade a mudanca & & & \\
Nao e reconhecido por testes verbais comuns & & & \\
Sinais de lateralidade e praxia alteradas & & & \\
\bottomrule\end{tabular}
\subsection*{Como investigar (tarefas sistematicas)}
\begin{itemize}
\item Observacao de habilidades visuoespaciais
\item Avaliacao psicopedagogica com testes verbais vs visuomotores
\item Entrevista com pais
\end{itemize}
\subsection*{Exclusoes e confundidores}
\begin{itemize}
\item TEA (pragmatica comprometida em TANV e TEA)
\item Disgrafia isolada
\item TDL verbal (perfil oposto)
\end{itemize}
\subsection*{Instrumentos validados para confirmacao (referencia -- aplicar somente por psicologo, conforme SATEPSI; nao reproduzidos aqui)}
\begin{itemize}
\item Avaliacao neuropsicologica (psicologo)
\item Avaliacao perceptomotora (terapeuta ocupacional)
\end{itemize}
\subsection*{Encaminhamento}
\begin{itemize}
\item Psicologo/neuropsicologo
\item Terapeuta ocupacional
\end{itemize}
\subsection*{Adaptacoes pedagogicas iniciais (sem rotulo)}
\begin{itemize}
\item Ensino explicito de estrategias visuoespaciais
\item Organizadores graficos para matematica
\item Antecipar mudancas com rotina visual
\item Ensino de habilidades sociais explicitas
\item Tempo extra em tarefas de copia
\end{itemize}
\subsection*{Bibliografia basica}
\begin{itemize}
\item Rourke (1995), NLD. Guilford.
\item Margolis (2010). NLD and school. Springer.
\end{itemize}
\avisolegal


\chapter{Deficiencia Auditiva (sinais escolares)}
\section*{Dominio: Saude Sensorial}
\section*{Ficha de rastreio sistematico -- confirmacao}
\subsection*{Sinais de alerta (marque: S = sim, AV = as vezes, N = nao)}
\begin{tabular}{@{}p{11.5cm}ccc@{}}\toprule
Sinal observado & S & AV & N \\\midrule
Nao responde a chamado em volume normal & & & \\
Aumenta o volume de TV/rádio & & & \\
Pede repeticoes com frequencia & & & \\
Fala com qualidade alterada (voz, articulacao) & & & \\
Dificuldade de compreender em ruido & & & \\
Olha atentamente os labios do interlocutor & & & \\
Historico de otites de repeticao & & & \\
Uso de aparelho/implante (verificar funcionamento) & & & \\
Atraso academico em leitura e escrita & & & \\
Sinais detectados em triagem auditiva escolar & & & \\
\bottomrule\end{tabular}
\subsection*{Como investigar (tarefas sistematicas)}
\begin{itemize}
\item Triagem auditiva formal (OTORRINO/fonoaudiologo)
\item Observacao de respostas a sons agudos e graves
\item Verificar funcionamento do dispositivo DI/ATS
\end{itemize}
\subsection*{Exclusoes e confundidores}
\begin{itemize}
\item Distracao/desatencao
\item TDL sem perda auditiva
\item Ambiente ruidoso
\end{itemize}
\subsection*{Instrumentos validados para confirmacao (referencia -- aplicar somente por psicologo, conforme SATEPSI; nao reproduzidos aqui)}
\begin{itemize}
\item Audiometria (OTORRINO/fonoaudiologo)
\item Avaliacao de linguagem se perda confirmada
\end{itemize}
\subsection*{Encaminhamento}
\begin{itemize}
\item OTORRINO + fonoaudiologo (intervencao precoce)
\item Professor de educacao especial/AEE
\end{itemize}
\subsection*{Adaptacoes pedagogicas iniciais (sem rotulo)}
\begin{itemize}
\item Assento estrategico (ouvido bom/visao do professor)
\item Falar de frente, com luz no rosto
\item Reduzir ruido de fundo
\item Uso de legendas/apoio visual
\item Checklist de funcionamento do ATS/diario
\end{itemize}
\subsection*{Bibliografia basica}
\begin{itemize}
\item Boothroyd (2004). Hearing loss. Plural.
\item Bevilacqua (2004). Deficiencia auditiva. Pulso.
\end{itemize}
\avisolegal


\chapter{Baixa Visao (sinais escolares)}
\section*{Dominio: Saude Sensorial}
\section*{Ficha de rastreio sistematico -- confirmacao}
\subsection*{Sinais de alerta (marque: S = sim, AV = as vezes, N = nao)}
\begin{tabular}{@{}p{11.5cm}ccc@{}}\toprule
Sinal observado & S & AV & N \\\midrule
Espreme os olhos ou aproxima muito o papel & & & \\
Queixa de dor de cabeca ou vista cansada & & & \\
Dificuldade de copiar do quadro & & & \\
Perde linha ao ler & & & \\
Inclina a cabeca para ver & & & \\
Sensibilidade a luz ou claridade excessiva & & & \\
Erros em tarefas de detalhe (labirintos) & & & \\
Esbarra ou desvia de obstaculos com frequencia & & & \\
Nao percebe objetos distantes & & & \\
Historico de uso de oculos (verificar atualizacao) & & & \\
\bottomrule\end{tabular}
\subsection*{Como investigar (tarefas sistematicas)}
\begin{itemize}
\item Triagem visual (ametropia) com profissional
\item Observacao de distancias de trabalho
\item Rever data da ultima consulta oftalmologica
\end{itemize}
\subsection*{Exclusoes e confundidores}
\begin{itemize}
\item Dificuldade atencional
\item Desinteresse
\item Problema motor (nao visual)
\end{itemize}
\subsection*{Instrumentos validados para confirmacao (referencia -- aplicar somente por psicologo, conforme SATEPSI; nao reproduzidos aqui)}
\begin{itemize}
\item Avaliacao oftalmologica (exame completo)
\item Testes de acuidade visual
\end{itemize}
\subsection*{Encaminhamento}
\begin{itemize}
\item Oftalmologista (consulta e correcao)
\item Professor de educacao especial/AEE se baixa visao confirmada
\end{itemize}
\subsection*{Adaptacoes pedagogicas iniciais (sem rotulo)}
\begin{itemize}
\item Assento proximo ao quadro com boa iluminacao
\item Material ampliado (fonte 16-18pt no livro)
\item Contraste alto no material
\item Uso de recursos de ampliacao (lupa, tela)
\item Reduzir brilho e ofuscamento na sala
\end{itemize}
\subsection*{Bibliografia basica}
\begin{itemize}
\item Veitzman (2007). Visao na escola. Santos.
\item Colaris (2020). Baixa visao. UFPR.
\end{itemize}
\avisolegal


\chapter{Comunicacao Social (TCS) pragmatica}
\section*{Dominio: Linguagem/Social}
\section*{Ficha de rastreio sistematico -- confirmacao}
\subsection*{Sinais de alerta (marque: S = sim, AV = as vezes, N = nao)}
\begin{tabular}{@{}p{11.5cm}ccc@{}}\toprule
Sinal observado & S & AV & N \\\midrule
Fala fluente mas socialmente desajeitada & & & \\
Dificuldade de interpretar figurativo, ironia e ambiguidade & & & \\
Troca de turnos pobres na conversa & & & \\
Discurso monologico sobre interesses & & & \\
Dificuldade de adequar registro ao interlocutor & & & \\
Problemas para fazer e manter amizades & & & \\
Sem restricoes de interesses ou comportamentos sensoriais (contraste com TEA) & & & \\
Compreensao literal excessiva & & & \\
Dificuldade de inferencia em textos & & & \\
Prejuizo funcional em interacoes & & & \\
\bottomrule\end{tabular}
\subsection*{Como investigar (tarefas sistematicas)}
\begin{itemize}
\item Observacao de conversacao em pares
\item Prova de compreensao pragmatica (figurativo, ironia)
\item Comparar com TEA (ausencia de repeticoes/restricoes)
\end{itemize}
\subsection*{Exclusoes e confundidores}
\begin{itemize}
\item TEA (tem restricoes e padroes sensoriais)
\item TDL (gramatica comprometida)
\item Timidez/ansiedade social
\end{itemize}
\subsection*{Instrumentos validados para confirmacao (referencia -- aplicar somente por psicologo, conforme SATEPSI; nao reproduzidos aqui)}
\begin{itemize}
\item Avaliacao fonoaudiologica de pragmatica
\item Avaliacao clinica (psicologo)
\end{itemize}
\subsection*{Encaminhamento}
\begin{itemize}
\item Fonoaudiologo (pragmatica)
\item Psicologo se comorbidade social/emocional
\end{itemize}
\subsection*{Adaptacoes pedagogicas iniciais (sem rotulo)}
\begin{itemize}
\item Ensinar explicitamente interpretacao de figurativo e ironia
\item Role-play de trocas sociais com feedback
\item Grupo de conversacao estruturado
\item Leitura de historias com perguntas inferenciais
\item Reforcar tentativas sociais adequadas
\end{itemize}
\subsection*{Bibliografia basica}
\begin{itemize}
\item DSM-5-TR (2023). Artmed.
\item Adams (2002). Social communication. Wiley.
\end{itemize}
\avisolegal



\chapter{Sintese, Fluxo e Orientacao de Professores}

\section{Folha de Sintese}
\begin{quote}
Preencher apos a Sondagem (Parte A) e as fichas de confirmacao (Parte B).
\end{quote}

\noindent Data: \rule{3cm}{0.3pt}\quad Aluno: \rule{6cm}{0.3pt}\quad Idade: \rule{1.5cm}{0.3pt}

\subsection*{Dominios com prioridade de investigacao (maior soma na Parte A)}
\begin{enumerate}[leftmargin=2.2em]
\item \rule{12cm}{0.3pt}
\item \rule{12cm}{0.3pt}
\item \rule{12cm}{0.3pt}
\end{enumerate}

\subsection*{Fichas de confirmacao aplicadas}
% PAGAR campo
\noindent Ficha(s): \rule{14cm}{0.3pt}

\subsection*{Hipoteses de trabalho (nunca diagnosticos)}
\noindent \rule{15cm}{0.3pt}\\ \rule{15cm}{0.3pt}

\subsection*{Encaminhamentos}
\begin{itemize}
\item Servico de saude / profissional: \rule{6cm}{0.3pt}\quad Data: \rule{2.5cm}{0.3pt}
\item AEE: \rule{10cm}{0.3pt}
\end{itemize}

\subsection*{Aviso de sigilo e uso restrito}
\avisolegal

\section{Fluxo de Encaminhamento}
\begin{enumerate}[leftmargin=2.2em]
\item \textbf{Sondagem Inicial} (Parte A), antes da primeira atividade.
\item Dominios com maior soma -> \textbf{fichas de confirmacao} (Parte B).
\item Ficha com sinais relevantes -> \textbf{encaminhamento} ao profissional habilitado.
\item Professor recebe \textbf{orientacoes de adaptacao} (nao rotulos).
\item Reavaliar apos 2-3 meses de intervencao; registrar progressos.
\end{enumerate}

\section{Orientacoes ao Professor (material do aluno permanece sem avaliacao clinica)}
\begin{itemize}
\item O aluno NAO deve ter contato com este volume.
\item O professor NAO aplica testes; ele observa, registra e adapta.
\item Adaptacoes sao planejamento pedagogico, nao e diagnostico.
\item Qualquer suspeita e comunicada ao profissional responsavel pelo volume.
\item Elogiar esforco e progresso; jamais usar fichas como rotulo ou punicao.
\end{itemize}

\avisolegal


\chapter*{Referencias de Instrumentos Citados}
\addcontentsline{toc}{chapter}{Referencias de Instrumentos Citados}

\section*{Como usar esta tabela}
Esta secao reune, por referencia (SATEPSI/CFP e literatura cientifica), os
instrumentos citados nas fichas da Parte B. Nenhum instrumento e reproduzido
neste volume: o profissional habilitado aplica a versao original adquirida junto
a editora detentora dos direitos. A indicacao de quem aplica segue o SATEPSI
(CFP) e as normas de cada editora. DOIs conferidos na data de producao desta
edicao; revisar antes de citar em laudos.

\section*{Tabela de Referencia}
\setlength{\tabcolsep}{4pt}
\begin{longtable}{@{}>{\raggedright\arraybackslash}p{3.0cm}>{\raggedright\arraybackslash}p{3.6cm}>{\raggedright\arraybackslash}p{1.8cm}>{\raggedright\arraybackslash}p{2.7cm}>{\raggedright\arraybackslash}p{4.1cm}@{}}
\toprule
Instrumento & O que avalia & Faixa & Quem aplica & Referencia / DOI \\
\midrule
\endfirsthead
\toprule
Instrumento & O que avalia & Faixa & Quem aplica & Referencia / DOI \\
\midrule
\endhead
TDE-II (Stein, Giacomoni e Fonseca) & Leitura, escrita e aritmetica escolar & 1.o ao 9.o ano & Psicopedagogo, neuropsicopedagogo, psicologo e fonoaudiologo (SATEPSI: nao privativo) & Vetor Editora. Unico instrumento psicopedagogico normalizado para o Brasil. \\
SRS-2 & Responsividade social (triagem TEA) & 2,5 anos a adulto (versoes) & Psicologo & Barbosa et al. (2015). J Bras Psiquiatr. DOI: 10.1590/\hspace{0pt}0047-2085000000083. Otoni et al. (2023). Psic.: Teor. e Pesq. DOI: 10.1590/\hspace{0pt}0102.3772e39nspe11.en. \\
ADI-R & Entrevista diagnostica TEA com responsaveis & Criancas e adolescentes & Psicologo/psiquiatra habilitado & Becker et al. (2012). Arq Neuro-Psiquiatr. DOI: 10.1590/\hspace{0pt}S0004-282X2012000300006. \\
ADOS-2 & Avaliacao observacional padronizada TEA & Criancas e adolescentes & Psicologo/psiquiatra com formacao especifica & Lord et al. (2012). WPS. Em processo de validacao no Brasil; usar normas internacionais com cautela. \\
CARS-BR & Gravidade de sintomas de TEA & 3 a 17 anos & Psicologo/psiquiatra & Traducao/validacao brasileira (LUME UFRGS); alfa de Cronbach 0,82; concordancia com ATA r = 0,89. \\
Mini-TEA & Rastreio de TEA & 2,5 a 12 anos & Profissional de saude habilitado & Escala brasileira recente; ver PubMed 38438070. \\
SNAP-IV & Sintomas de TDAH e TOD (relato de pais e professores) & 4 a 16 anos & Triagem (uso profissional; nao diagnostica) & Mattos et al. (2006). Rev Psiquiatr RS. DOI: 10.1590/\hspace{0pt}S0101-81082006000300008. Costa et al. (2019). J Pediatr (Rio J). DOI: 10.1016/\hspace{0pt}j.jped.2018.06.014. Dominio publico. \\
CBCL/TRF & Problemas de comportamento (internalizantes e externalizantes) & 6 a 18 anos & Psicologo & Achenbach (ASEBA); adaptacao brasileira publicada em periodicos indexados. \\
SDQ & Triagem de saude mental (forcas e dificuldades) & 4 a 17 anos & Uso profissional em triagem & Goodman; versao brasileira validada. \\
ABFW & Linguagem infantil: fonologia, vocabulario, fluencia, pragmatica & Infantil (provas por faixa etaria) & Fonoaudiologo & Andrade, Befi-Lopes, Fernandes e Wertzner. Pro-Fono. \\
PROLEC & Processos de leitura (letras, lexico, sintaxe, semantica) & 2.o ao 5.o ano & Fonoaudiologo, psicologo e professor (adaptacao brasileira) & Cuetos, Rodriguez e Ruano; adapt. Capellini, Oliveira e Cuetos. Hogrefe. Nota: PROLEC-T tem amostra limitada no Brasil (Pinheiro, 2017); uso informal. \\
MABC-2 & Coordenacao motora (destreza manual, mirar e pegar, equilibrio) & 3 a 16 anos (BI2: 7 a 10) & Psicologo, terapeuta ocupacional, fisioterapeuta, educacao fisica, pediatra & Henderson, Sugden e Barnett. Pearson. \\
CDI & Sintomas depressivos (autorrelato) & 7 a 17 anos & Psicologo & Kovacs (1983); adaptacao brasileira Gouveia et al. (1995); versao reduzida em Cruvinel e Boruchovitch (2008). DOI: 10.1590/\hspace{0pt}S1414-98932008000300011. \\
CABI & Comportamento e psicopatologia (relato parental) & 6 a 18 anos & Psicologo (triagem) & Adaptacao brasileira publicada em 2023; ver PMC10373135. Gratuito. \\
WISC-V & Inteligencia e processamento cognitivo & 6 a 16 anos & Psicologo (SATEPSI privativo) & Wechsler. Pearson. \\
\bottomrule
\end{longtable}

\section*{Avisos de rigor}
\begin{itemize}
\item O ADOS-2 ainda nao possui normas brasileiras completas; sua mencao e
referencial, para uso por profissional com formacao especifica.
\item O PROLEC-T (prova de compreensao de texto) apresenta amostra limitada no
Brasil e deve ser usado apenas como avaliacao informal (Pinheiro, 2017).
\item Instrumentos privativos de psicologo (ex.: WISC-V, CDI, CBCL) seguem a
Resolucao CFP n. 009/2018 e o SATEPSI; psicopedagogos e neuropsicopedagogos
aplicam os instrumentos nao privativos (ex.: TDE-II) e entrevistas/observacao.
\item Este volume nao reproduz itens, normas ou materiais de nenhum instrumento;
as fichas da Parte B sao observacao pedagogica propria e rastreio sistematico.
\end{itemize}

\avisolegal


\backmatter
\end{document}
